\begin{abstract}
A representação da madeira por classes de resistência simplifica o dimensionamento, mas sua influência sobre o desempenho e o consumo de material depende das propriedades substituídas e das verificações que condicionam o projeto. O enquadramento é indexado pela resistência característica à compressão paralela às fibras, enquanto a rigidez acompanha por associação, sem participar do critério. Este estudo quantifica quanto a substituição das propriedades medidas de uma espécie pelas propriedades da sua classe altera o volume de madeira e o atendimento ao limite de serviço em superestruturas de pontes de madeira roliça para estradas vicinais. A base experimental é o conjunto de 40 espécies folhosas nativas caracterizadas por Dias e Rocco Lahr (2004), cujas resistências características publicadas permitem reproduzir o enquadramento pelo mesmo critério aplicado em projeto. A comparação é conduzida sobre o mesmo modelo estrutural e carregamento, substituindo o conjunto completo de propriedades da espécie pelos valores da classe, o que reproduz a decisão de projeto usual. Cada cenário é redimensionado por otimização multiobjetivo sob critérios de aceitação idênticos, com tolerância dimensional de 5\% no diâmetro avaliada no pior caso, e repetido com sementes pareadas, de modo a separar efeito de material da dispersão entre execuções do otimizador. Os projetos obtidos pela classe são então reavaliados com a rigidez medida, sem reotimização, o que distingue diferença de consumo de material de excedência do limite de deslocamento. A campanha cobre vãos de 3 a 10~m. \tofill{[Após as análises: informar erro de previsão de flecha, diferença de volume entre projeto pela espécie e pela classe, frequência de excedência do limite de serviço na reavaliação e sensibilidade ao vão e à escolha do módulo.]}
\end{abstract}

\textbf{Palavras-chave}: pontes de madeira; madeiras tropicais; classes de resistência; rigidez; estado limite de serviço; consumo de material.
